% NAL 1644, batch 1 — Gallica views 1–20.
% Pass: Opus 5.5 (claude-opus-5-5), 2026-09-23 — first pass, unchecked against
% the views by a human.
%
% What the views are. Black-and-white scans, portrait, about 3748 × 5622, one
% page per view. The leaves are of several sizes and sit on a white ground;
% smaller leaves let the edges of the leaves beneath or beside them show
% (views 2–4, 14, 18, 20), and show-through from the other side is strong on
% every written leaf. Each view was first seen as a thumbnail, then read as
% four full-width bands at 2400 px and, for the text block, as native-width
% strips of about 470 px height, with tight crops for struck words, the Greek,
% the figures and the numerals. Read from the local mirror only. \page{N} is
% always the Gallica view.
%
% Dating. The catalogue gives none; \dating{} is left empty, and this pass
% infers no date. The Elzevir edition of 1646 named in the notes of views 1
% and 5 is a later reader's reference, not a date of the leaves.
%
% What the twenty views hold.
%
%   v. 1      A paper wrapper (the front cover), torn in its lower half, with
%             the title « Francisci Vietæ opera nonnulla » in a large old
%             hand, and below, in a small later hand, a reference to « pag. 84
%             de l'edition d'Elzevir 1646 » and a two-line table: « de
%             æquationum recognitione 1. / De emendatione æquationum 32 ».
%             Shelfmarks and stamps. Transcribed, because it says what the
%             leaves were once thought to be.
%   v. 2–3    PIECE 1. Two small leaves lettered « A » and « B » in ink at the
%             top right, in a compact hand different from the treatise's, in
%             Latin. v. 2: the « duplicata hypostasis » and the « climactica
%             paraplerosis » (reduction of E cc + Z s. in E c. æq. D ss. and
%             of E cc + Z ppl. in E q. æq. D ss. to proportionals), then
%             « Facillima resolutio æquationum cubicarum affectarum »,
%             reducing them to cubics « a radice simplici ». v. 3: a theorem
%             on sums and differences of alternate terms in a series of
%             continued proportionals, even or odd in number, « quod methodo
%             Analyticâ facile patebit »; the rest of the leaf is blank. A
%             margin note on v. 2, in blacker ink and another hand, reads
%             « Vid. fol. 2 ».
%   v. 4      A third such leaf, lettered « C », blank. Not transcribed.
%   v. 5–20   PIECE 2. « Francisci Vietæ De Recognitione Æquationum
%             tractatus. », so headed on the leaf, in one regular, posed
%             cursive in pale ink, Latin with Greek words in Greek letters. It
%             runs without a break from v. 5 to v. 20 and does not end there:
%               Cap. 1 (v. 5) De dignoscenda æquationum constitutione ex
%                 Zetesi, plasmate et Syncrisi;
%               Cap. II (v. 6) De Zetesi; similar equations defined;
%               Cap. III (v. 6–8) quadratics from the zetetics, three
%                 theorems (καταφατική, ἀποφατική, ἀμφίβολος), numerical
%                 examples on the series 8, 12, 18;
%               Cap. IIII (v. 8–10) cubics affected under the side, three
%                 theorems, series 8, 12, 18, 27 and a « double » second term
%                 (series L.59319, 195, L.24375, 125 and L.59319, 78, L.624, 8);
%               Cap. V (v. 10–13) cubics affected under the square, three
%                 theorems;
%               Cap. VI (v. 13–17) « Alia insuper Constitutio Cuborum sub
%                 latere adfectorum », with the condition on the quadruple
%                 cube of a third of the plane coefficient, two theorems,
%                 their geometric restatement « Aliter », a third theorem for
%                 the irreducible case, and « Aliter, tertium theorema »
%                 through two right triangles whose angles are as 3 to 1,
%                 example 1C − 300N æq. 432 (roots 18 and 9 ± L57);
%               Cap. VII (v. 18–20) « De Generali methodo transmutandarum
%                 æquationum »: transmutation by alteration of the root (seven
%                 modes and composite ones) and with the root unaltered
%                 (climactic ascent and descent, regular and irregular);
%               Cap. VIII (v. 20) « Singularia de plasmate », begun.
%             The text of v. 20 stops without punctuation on « si quo id
%             liceat compendio »: it runs on into view 21 (the verso, in
%             batch 2). The same holds for the page breaks at v. 6→7, 8→9,
%             10→11, 12→13, 14→16 (v. 15 is an inserted slip) and 16→17.
%
% A second hand. A rapid, smaller hand in blacker ink annotates the treatise
% from v. 13 on: an addition at the foot of v. 13 (« adscripta conditio hic
% patet », with « vide Zet. 20 lib. 2. » in the margin); two margin notes on
% v. 14; the whole of the inserted slip v. 15, headed « vel sic » and marked
% « 5 bis », a proof of that condition; the left margin and foot of v. 16
% (a first, struck-through draft of the same proof, and a justification of
% Theorem III); on v. 17–18 a figure (a circle on the diameter AB, the arc FB
% trisected at G, K) and a long proof of « Aliter, tertium theorema » running
% down both margins. This hand writes « è demonstratis a nobis » and
% « demonstratum a nobis ad sectiones angulares » (v. 17–18): its writer
% speaks of his own work on angular sections. The leaves do not name him, and
% this pass attributes neither this hand nor the treatise's hand to anyone;
% the only name on the leaves is the heading's and the cover's « Vieta ». A
% third, small and pale hand has left three short notes keyed by « + » in the
% left margin of v. 20. Heavy black ink, perhaps a later reader's, repasses a
% few letters of the text and writes a capital « D » in the margin of v. 12
% and a circled « D » at the foot of v. 16 and the head of v. 17, which there
% serves as a cross-reference between the two parts of the struck proof.
%
% Notation, kept as written. The unknown is A (and E after a transformation);
% the given magnitudes are consonants, named with their dimension: « B
% planum », « D solidum », « Z quadratum », abbreviated « p. », « pl. »,
% « s. », « pp. », « ppl. », « ss. » in the small leaves; « in » for a product;
% « æquatur », « æquabitur », or « æq. » for equality — there is no sign of
% equality. In the numerical examples the cossic letters N (the root), Q
% (square, written like an underlined 2), C (cube) are used, and « L » for the
% side (root): « 1C − 64N. æquetur 960 », « L. 144 ». The second hand writes
% the powers A q, A c, A qq, A qc, A cc (its c looks like an r) and puts a
% numerical coefficient after the term (« B c in A c 26 »). Sums written in
% columns are set as arrays, as on the leaf. Roman numerals I–IIII stand over
% the terms of each numerical series and are set with \overset. The long s is
% set s. One glyph has no LaTeX equivalent: the cross-reference sign of the
% second hand, a saltire with a small circle above followed by « 4 » (v. 15,
% 16), set $\overset{\circ}{x}\,4$.
%
% Numbers on the leaves.
%   - A heavy ink figure at the top right of the treatise's rectos: 1 (v. 5),
%     2 (v. 7), 3 (v. 9), 4 (v. 11), 5 (v. 13), « 5 bis » on the inserted slip
%     (v. 15, in the second hand's ink), 6 (v. 16), 7 (v. 18), 8 (v. 20). One
%     per leaf, rectos only, starting at 1 with the treatise and not counting
%     the lettered leaves A–C: a numbering of the treatise's leaves, not the
%     library's foliation. Figures « 9 » and « 11 » show on the edges of
%     further leaves projecting at the right of v. 20, apparently of the same
%     series. Nothing here is a pale pencil foliation, so no \folio{} is
%     written anywhere in the batch.
%   - Letters A, B, C at the top right of the small leaves (v. 2, 3, 4).
%   - Shelfmarks and stamps: « Libri 1201 », « Lat. nouv. acq. 1644 »,
%     « Nouv. acq. lat. 1644 », « Lat. n. a. 1644 », « R.c. 8070 (73) », the
%     oval stamp « ACQ. 8070 (LIBRI) », the round stamp « Bibliothèque
%     nationale Manuscrits » (v. 1, 2, 5, and a pale round stamp on v. 15 and
%     v. 16). A pencil « 8 » near the lower left edge of the cover (v. 1).
%   - References: « pag. 84 … p 158 » of the Elzevir 1646 edition (v. 1),
%     « page 84 de l'edition d'Elzevir 1646 in folio » (v. 5); the table « 1 »
%     and « 32 » on the cover; « Vid. fol. 2 » (v. 2); « vide Zet. 20 lib. 2 »
%     (v. 13). None of these was checked against any edition.
%
% Not settled. Whether the lettered leaves A–C and the treatise belong to one
% dossier; whose the three hands are; the Greek words, several of which are
% offered with doubt; the long margins of v. 16–18, read with many \ill{} and
% \uncertain{}.

% Coordinator's reconciliation (2026-09-23), after all nine batches:
%   One ink series of leaf numbers runs top right of the rectos through the
%   whole volume, with no gap: 1–8 and « 5 bis » (v. 5–20), 9–17 (v. 22–39),
%   18–27 (v. 41–59), 28–37 (v. 61–79), 38–47 and « 46 bis » (v. 81–100),
%   « 47 bis » and 48–56 (v. 102–120), 57–66 (v. 122–140), 67–78
%   (v. 142–160), 79 (v. 162). It crosses the three pieces of the volume, and
%   the writer cites it himself (« fol. 18. », v. 43; « pag. 48 », v. 100), so
%   it is the writer's or copyist's count of leaves, not the library's. Being
%   ink, it gets no \folio{}, in any batch.
%   The pieces: « Francisci Vietæ De Recognitione Æquationum tractatus »
%   (heading on v. 5), Cap. 1–XXI, v. 5–68; « Francisci Vietæ. De æquationum
%   Emendatione tractatus secundus » (heading on v. 69), Cap. I–XIIII, ending
%   « Finis » on v. 141; then astronomy without a name, v. 142–163 (lunar
%   prostaphaereses, Tycho's lunar hypothesis restated, « Harmonici
%   Ptolemaici Constructio Caput XII », Mercury). Two doubtful notes may mark
%   the treatises as a copy: « transcripta ex exemplari » (v. 100) and
%   « deerant hæc in exemplari » (v. 114). The name on the leaves is in the
%   headings; who wrote the hands is not settled.

\documentclass[11pt,a4paper]{article}
% The preamble shared by every transcription of the mathematicians' manuscripts in Gallica.
%
% It rests on one idea: **uncertainty compounds, it does not resolve.** A
% transcription that smooths over a doubtful word has destroyed the only thing
% it was made for — a reader must be able to tell, at every point, what was on
% the page and what was supplied. The apparatus macros below are therefore the
% critical apparatus, not decoration.
%
% The same commands are recognised by `scripts/render.mjs`, which produces the
% site's reading view, and by `scripts/tei.mjs`. A macro added here must be
% added there too, or rendering fails loudly — which is the wanted behaviour.
%
% Comments here are English, like the rest of the repository. The documents
% this preamble sets are French, and that is deliberate: see the skills.

\ProvidesPackage{germain}[2026/09/15 Transcriptions of mathematicians' manuscripts in Gallica]

\usepackage{amsmath,amssymb,amsthm}
\usepackage{mathrsfs}
% Kept from the parent preamble so the renderer's diagram support stays
% available; these manuscripts draw no commutative diagrams, and nothing here uses it.
\usepackage{tikz-cd}
\usepackage[english,french]{babel}
\usepackage{fontspec}

% --- Greek ------------------------------------------------------------------
% The text face stays fontspec's default, Latin Modern, which has no Greek:
% XeTeX drops every glyph it lacks with a line in the log and nothing on the
% page, and the Greek words the manuscripts quote vanished from the PDFs. So
% the Greek and Greek Extended blocks get a XeTeX character class of their own,
% and entering or leaving a run of it switches the family to CMU Serif — the
% same Computer Modern design, with polytonic Greek — and back. Only the
% family changes: a Greek word in \textit or \textbf keeps its shape and series.
% The transitions to and from every other class carry the tokens that class
% has towards an ordinary letter, so French spacing before « ; » or inside
% « » still holds after a Greek word. No grouping, so no brace or \endgroup
% can come out unbalanced; maths is left alone, since XeTeX inserts nothing
% there. Kept identical in `exercices.sty`.
\makeatletter
\newfontfamily\germainGreekfont{cmun}[
  Extension=.otf, NFSSFamily=germaingreek,
  UprightFont=*rm, ItalicFont=*ti, SlantedFont=*sl,
  BoldFont=*bx, BoldItalicFont=*bi, BoldSlantedFont=*bl]
\newXeTeXintercharclass\germain@greek
\def\germain@greekrange#1#2{%
  \@tempcnta=#1\relax
  \loop
    \XeTeXcharclass\@tempcnta=\germain@greek
    \ifnum\@tempcnta<#2\advance\@tempcnta\@ne\repeat}
\germain@greekrange{"0370}{"03FF}
\germain@greekrange{"1F00}{"1FFF}
\def\germain@greekprev{\rmdefault}
\def\germain@greekin{%
  \edef\germain@greekprev{\f@family}\fontfamily{germaingreek}\selectfont}
\def\germain@greekout{\fontfamily{\germain@greekprev}\selectfont}
\def\germain@greekjoin{%
  \XeTeXinterchartokenstate=\@ne
  \@tempcnta=\z@
  \loop
    \ifnum\@tempcnta=\germain@greek\else
      \XeTeXinterchartoks\germain@greek\@tempcnta=\expandafter{%
        \expandafter\germain@greekout\the\XeTeXinterchartoks\z@\@tempcnta}%
      \XeTeXinterchartoks\@tempcnta\germain@greek=\expandafter{%
        \the\XeTeXinterchartoks\@tempcnta\z@\germain@greekin}%
    \fi
    \ifnum\@tempcnta<4095\advance\@tempcnta\@ne\repeat}
% babel-french sets its own classes, and saves and restores the interchar
% state, each time a language is selected: join again after it does.
\addto\extrasfrench{\germain@greekjoin}
\addto\extrasenglish{\germain@greekjoin}
\AtBeginDocument{\germain@greekjoin}
\makeatother

\usepackage{geometry}
\geometry{a4paper,margin=25mm,marginparwidth=28mm}
\usepackage{marginnote}
\usepackage{xcolor}
\usepackage[normalem]{ulem}
\setlength{\emergencystretch}{3em}
\usepackage{hyperref}
\hypersetup{colorlinks=true,linkcolor=black,urlcolor=black,pdfborder={0 0 0}}

% --- Batch metadata ---------------------------------------------------------
% Taken as they stand by the rendering script, which shows them at the head of
% the reading view. They fix what the file claims to cover. The macro names are
% the parent project's, so that the scripts stay shared; \folder here means the
% volume's slug — `fr-9115` — and \pages counts Gallica views.

\newcommand{\folder}[1]{\gdef\grothFolder{#1}}
\newcommand{\batch}[1]{\gdef\grothBatch{#1}}
\newcommand{\pages}[2]{\gdef\grothFirst{#1}\gdef\grothLast{#2}}
\newcommand{\foldertitle}[1]{\gdef\grothFoldertitle{#1}}

% The holder's own shelfmark, printed as they write it: « Français 9115 ».
\newcommand{\shelfmark}[1]{\gdef\grothShelfmark{#1}}

% The dating, copied verbatim from the holder's catalogue — never inferred
% here. « XIXe siècle » is the BnF's; a pass that dates a leaf from its content
% says so in a \note{}, as its own inference.
\newcommand{\dating}[1]{\gdef\grothDating{#1}}

% The status notice, printed in the title block and repeated as a diagonal
% watermark on every page of the PDF. The manuscripts are in the public
% domain; what this line declares is not a right but a quality — an
% unverified first pass by a machine. The skills write the line; a file
% without it gets no watermark, so leaving it out is visible in the output.
\newcommand{\watermark}[1]{\gdef\grothWatermark{#1}}

\gdef\grothFolder{?}\gdef\grothBatch{}\gdef\grothFirst{?}\gdef\grothLast{?}%
\gdef\grothFoldertitle{}\gdef\grothShelfmark{}\gdef\grothDating{}\gdef\grothWatermark{}

% --- The résumé -------------------------------------------------------------
\newenvironment{resume}
  {\small\setlength{\parskip}{0.45em}}
  {\par\medskip\hrule\medskip}

% --- Keywords ---------------------------------------------------------------
% In English, closing the modernised reading's résumé: the modern vocabulary
% under which to search for what the leaves construct. The manifest extracts
% them to tag the volume — this line is the single source of the tags.
\newcommand{\keywords}[1]{\par\smallskip\noindent
  {\footnotesize\textsc{Keywords} --- \textit{#1}}\par}

% --- The critical apparatus -------------------------------------------------

% The Gallica view number, in the margin. This is the anchor that lets a
% disagreement over a formula be settled by looking at *one* image rather than
% a volume of 750. The site uses it to turn the facsimile as the transcription
% is scrolled. A view is one scanned image: usually one side of a leaf.
\newcommand{\page}[1]{%
  \marginnote{\footnotesize\color{gray!70}\textsf{v.\,#1}}%
  \label{page:#1}\ignorespaces}

% The BnF's pencilled foliation, where the leaf carries it and a pass can read
% it: « 198r ». This is what the literature cites — « ff. 198r–208v » — and
% the one line that turns a citation into a view. Recorded, never inferred:
% a folio a pass cannot read is not written.
\newcommand{\folio}[1]{%
  \marginnote{\footnotesize\color{gray!50}\textsf{f.\,#1}}\ignorespaces}

% The modernised reading's coarser counterpart to \page{}: which views a
% section is drawn from.
\newcommand{\pagerange}[2]{%
  \marginnote{\footnotesize\color{gray!70}\textsf{v.\,#1--#2}}%
  \ignorespaces}

% Illegible: never guessed.
\newcommand{\ill}{\textcolor{red!60!black}{[\,\ldots\,]}}

% A doubtful reading: the word is given, and flagged as uncertain.
\newcommand{\uncertain}[1]{\uline{#1}}

% An editorial addition — a missing letter, an implied word.
\newcommand{\add}[1]{\textcolor{blue!50!black}{[#1]}}

% Struck out by the writer. What was crossed out is part of the document.
\newcommand{\struck}[1]{\sout{#1}}

% The transcriber's note, on the state of the page or the mathematics.
\newcommand{\note}[1]{\par\smallskip{\small\color{gray!60!black}\textit{[#1]}}\par\smallskip}

% A marginal note of hers, distinct from the body of the text.
\newcommand{\marginal}[1]{\par\smallskip\noindent\hspace{1em}%
  {\small\color{gray!60!black}#1}\par\smallskip}

% --- Title ------------------------------------------------------------------
% The title block is French, like the documents themselves.

\AddToHook{shipout/background}{%
  \ifx\grothWatermark\empty\else
    \begin{tikzpicture}[remember picture, overlay]
      \node[rotate=52, scale=3.4, align=center, text=gray!18,
            font=\sffamily\bfseries]
        at (current page.center) {\grothWatermark};
    \end{tikzpicture}%
  \fi}

\AtBeginDocument{%
  \begin{center}
    {\large\bfseries Manuscrits de mathématiciens (Gallica) ---
      \ifx\grothShelfmark\empty\grothFolder\else\grothShelfmark\fi}\\[2pt]
    {\itshape\grothFoldertitle}\\[4pt]
    {\small vues \grothFirst--\grothLast
      \ifx\grothBatch\empty\else\ (lot \grothBatch)\fi}\\[2pt]
    \ifx\grothDating\empty\else
      {\small\color{gray!60!black}Datation du catalogue : \grothDating}\\[2pt]
    \fi
    \ifx\grothWatermark\empty\else
      {\small\bfseries\color{red!45!gray}\grothWatermark}\\[2pt]
    \fi
    {\footnotesize\color{gray!60!black}Transcription établie d'après les images IIIF de Gallica.
     Source gallica.bnf.fr / Bibliothèque nationale de France.\\
     Les lectures incertaines sont soulignées, les illisibles notées [\,\ldots\,],
     les ajouts entre crochets ; \textsf{v.} numérote les vues de Gallica,
     \textsf{f.} la foliotation au crayon de la BnF.}
  \end{center}
  \vspace{1em}}

\endinput


\folder{nal-1644}
\shelfmark{NAL 1644}
\batch{1}
\pages{1}{20}
\dating{}
\watermark{Lecture automatique\\non vérifiée}
\foldertitle{Viète (François). Papiers divers. NAL 1644}

\begin{document}

\page{1}
\note{Plat de couverture en papier fort, très taché, déchiré et lacunaire dans sa moitié inférieure, collé sur un feuillet de montage plus grand dont le bord blanc dépasse en bas. En haut à gauche, souligné d'une accolade, « Libri 1201 » ; en haut à droite, souligné d'une accolade, « Lat. nouv. acq. 1644 » ; au milieu, une grande lettre « G » à l'encre, et au-dessus, plus pâle, des traits qui paraissent un « G » plus ancien et une autre lettre, \ill{}. Au milieu de la moitié inférieure, d'une autre main, « Nouv. acq. lat. 1644 ». Cachet rond « Bibliothèque nationale Manuscrits » aux initiales « R.F. » ; en bas, le cachet ovale « ACQ. 8070 (LIBRI) » et, souligné d'une accolade, « R.c. 8070 (73) ». Un « 8 » au crayon clair près du bord gauche, en bas.}

Francisci Vietæ opera nonnulla

imprimés, \uncertain{pag. 84} de l'edition d'Elzevir 1646. in folio \uncertain{p} 158.

\uncertain{de} æquationum recognitione 1.

De \uncertain{emendatione} æquationum 32

\note{Le titre est d'une grande écriture cursive, d'une encre plus ancienne que les cotes. La ligne en français qui suit est d'une main plus petite, qui renvoie à l'édition d'Elzevir de 1646 ; le « 158 » est écrit sous « pag. 84 », précédé d'un trait descendant qui peut être un « p ». Les deux lignes de droite, de la même main que la note française semble-t-il, forment une petite table : le « 1 » et le « 32 » paraissent renvoyer à des pages du volume ; un trait clair (crayon ?) en forme de crochet est posé sur le « 32 ». Le deuxième titre est lu avec doute : les jambages laissent voir « em…ntione ».}

\page{2}
\note{Feuillet de papier plus petit que la page, monté ; au-dessus de son bord supérieur paraît la bande d'un autre papier portant à droite une lettre « C » à l'encre (la même bande, avec le même « C », se voit aux vues 3 et 4). À droite de la première ligne, une lettre « A » isolée, de la main du texte semble-t-il. Dans la marge de gauche, d'une encre plus noire et d'une autre main, « Vid. fol. 2 », souligné. Au-dessus de la première ligne, souligné d'une accolade, « Libri 1201 » ; à droite, souligné d'une accolade, « Lat. nouv. acq. 1644 ». Au pied, « R.c. 8070 (73) » souligné d'une accolade, le cachet ovale « ACQ. 8070 (LIBRI) » et le cachet rond « Bibliothèque nationale Manuscrits » aux initiales « R.F. ». Aucune foliation au crayon n'est visible.}

Opus \uncertain{duplicatæ} hypostaseos ad mechanicam Geometriam sit facile reducere.
\[ \begin{array}{l} E\,cc \\ + Z\,s.\ \text{in}\ E\,c.\ \text{æq.}\ D\,ss. \end{array} \]

Si fiat $D\,s.$ ad $Z\,s.$ ita quævis recta $D$ ad $B$. Et facta $D$ media et $B$ differentia extremarum, inveniantur extremæ quarum minor sit $F$. Denique ut \struck{$D\,s.$ ad $F\,s.$} $D$ ad $F$ ita fiat $D\,s.$ ad $E\,s.$ \uncertain{et est} inter $D$ et $F$. Inveniantur duæ mediæ continue proportionales quarum secunda erit ipsa $E$ quæsita.
\note{Le $D$ de « recta $D$ ad $B$ » est repassé à l'encre, sur une autre lettre. Les deux mots lus « et est » sont liés et incertains.}

\subsection*{Climactica paraplerosis.}

\[ \begin{array}{l} E\,cc. \\ + Z\,ppl.\ \text{in}\ E\,q.\ \text{æq.}\ D\,ss. \end{array} \]
\marginal{in 2}
\note{Dans la marge de gauche, en regard de cette équation, « in 2 » ; le chiffre a la forme que la main donne aussi à la lettre $Z$, et l'on peut lire « in Z ».}

fiat ut $Z\,ppl.$ ad $D\,pp.$ \struck{\uncertain{D est}} constituta serie continue proportionalium quarum $Z$ prima, $D$ secunda, fiat ut $Z$ ad quartam ita $D\,q.$ ad $C\,q.$ \uncertain{\&} ut $Z\,p.$ ad $C$ plano ita $Z$ ad $F$. Dataq; $Z$ prima \& serie quatuor continue proportionalium, $F$ composita ex secunda et quarta, inveniantur proportionales, et sit ut $Z$ ad secundam ita $Z$ pla. ad aliud quod æquale erit ipsi $E\,q.$ quæsito.
\note{« constituta » est écrit sur une première forme du mot ; les deux lettres lues « D est » qui précèdent sont barrées d'un trait qui les lie au mot suivant. La lettre lue $C$ (« $C\,q.$ », « $C$ plano ») pourrait être un $G$.}

\subsection*{Facillima resolutio æquationum cubicarum affectarum. \& opere duplicatæ hypostaseos \uncertain{restitutarum}.}
\note{Le titre est précédé, à gauche, de deux traits de plume obliques.}

\struck{\ill{}} Quemadmodum æquationes cubicæ quorumcunq; tandem modo affectæ ad quadraticas deprimi possunt \& radice simplici. Proponatur. $E$ \uncertain{plani} cubus plus minusve $Z\,s.$ in $E\,c.$ æq. \struck{$Z$} $B\,ss.$ fingatur $\dfrac{E\,c.}{Z}$ æquari $A\,q.$ atq; ita ordinetur æquatio ut \struck{\ill{}} \struck{\ill{}} æquatio namq; orta sic se habebit $A\,c. + Z\,q.$ in $A$ æq. $\dfrac{B\,ss.}{Z\,c.}$
\note{Le premier mot du paragraphe, avant « Quemadmodum », est surchargé et illisible. Le $B$ de « $B\,ss.$ » est écrit fortement sur un $Z$ barré ; la finale de « $B\,ss.$ » compte un jambage de plus, et l'on pourrait lire « $B\,sss.$ ». Les deux mots barrés de la ligne suivante sont couverts d'un trait épais.}

Sit \& opere climactices parapleroseos si reducatur in æquationem \uncertain{similem} $E$ plani cubus $+$ vel \struck{\ill{}} $-$ $Z\,ppl.$ in $E\,q.$ æq. $B\,ss.$ supponatur $\dfrac{E\,q.}{Z}$ esse $A$ et reducetur in æquationem cubicam a radice simplici \struck{\ill{}} \uncertain{ordinatæ}
\[ \begin{array}{l} A\,c. \\ + Z\,q.\ \text{in}\ A\ \text{æq.}\ \dfrac{B\,ss.}{Z\,c.} \end{array} \]
\note{En marge de gauche, devant « et reducetur », un signe en forme de croix double, comme un renvoi. Le mot barré après « simplici » est souligné de deux traits et barré ; on devine « quadr… ». Le signe d'égalité est partout l'abréviation « æq. » ou un grand signe bouclé qu'on transcrit « æq. ».}

atq; ita æquationes omnes a radice planâ solidave aut \ill{} climactica \uncertain{restituuntur} ad æquationem a radice simplici.

\page{3}
\note{Même disposition qu'à la vue 2 : un petit feuillet monté, au-dessus duquel dépasse la bande portant « C ». Le petit feuillet porte à droite, au-dessus du titre, une lettre « B » à l'encre. Au bord gauche de la vue paraissent, coupés, les mots « Vid. » et « fol. » de la marge de la vue 2 : le feuillet A déborde sous celui-ci. Le texte n'occupe que le tiers supérieur ; le reste de la page est blanc. Aucun numéro de folio.}

\subsection*{Theorema}

Si fuerit series continue proportionalium numero parium, differentia postremarum cum \add{alterne sumptarum} accepta prima, minus differentia cum alterne sumptarum æquatur extremæ minori plus eadem maiori.

at si proportionales fuerint numero impares, differentia cum alterne sumptarum minus differentia cum postremarum \add{alterne sumptarum} accepta prima, æquatur extremæ minori, plus eadem maiori. Debet autem in serie proposita intelligi prima minor. quod methodo Analyticâ facile patebit.
\note{Les deux « alterne sumptarum » en \emph{add} sont des ajouts interlinéaires de la main du texte, écrits au-dessus de « accepta », avec un signe d'insertion dans la première phrase. Le mot transcrit « cum » est partout abrégé en trois jambages et une boucle. Dans la première ligne, « series » est suivi d'un trait qui paraît une lettre reprise.}

\page{5}
\note{Nouveau feuillet, d'une autre main que les vues 2–3 : une cursive posée, régulière, d'encre pâle. En haut à droite, à l'encre noire et épaisse, le chiffre « 1 », d'un trait plus fort que le texte. À gauche du titre, souligné d'une accolade, « Libri 1201 » ; après le nom, souligné d'une accolade, « Lat. n. a. 1644 ». Cachet rond « Bibliothèque nationale Manuscrits » aux initiales « R.F. » au milieu de la page. Dans la marge de gauche, d'une petite main à l'encre noire, la même semble-t-il que celle de la note de la couverture : « page 84 de l'edition d'Elzevir 1646 in folio. » ; plus bas, une petite lettre « m. » isolée, et « R.c. 8070 (73) » souligné d'une accolade. L'écriture du verso transparaît dans la marge et entre les lignes.}

\section*{Francisci Vietæ De Recognitione Æquationum tractatus.}

\subsection*{De dignoscenda æquationum constitutione ex Zetesi, plasmate et Syncrisi. Cap. 1}

Generalem et generaliter traditam de numerosa potestatum resolutione doctrinam informat ac perficit tractatus de recognitione æquationum. Præparatione enim indigent æquationes sæpenumero antequam faciliter explicentur. ac præsertim cum potestates de homogeneis magnitudinibus negantur vel ita mixtim homogeneæ magnitudines de potestatibus negantur et affirmantur et affectiones negatæ affirmatis præpollent ac demq; quotquot æquationes fractionibus vel asymmetriis explicentur.

In Geometricis quidem \uncertain{ari}\ill{} fractione vel asymmetria non \uncertain{solet} æquationibus \uncertain{offerre}, quo \uncertain{minus} εὐμηχάνως explicentur. Sicuti neq; vitio negationis. Est enim certum \uncertain{semper} subiectum sub quo operatur Geometra, at abest πολυ\ill{}, et quo elatior est potestas, affectionisq; gradus, eo maior se prodit in explicando problemate \uncertain{ἀρρηξία} ἢ ἀλογία.
\note{Le cachet couvre la fin du mot après « quidem » et le début de « fractione ». Après « non », le mot lu « solet » a sa finale repassée d'un trait. Les trois mots grecs sont en caractères grecs dans la ligne ; le premier, après « abest », commence par « πολυ » et sa fin n'est pas lue ; celui qui précède « ἢ ἀλογία » est lu avec doute.}

\uncertain{Ergo} vero æquationi quæ proposita est agnita constitutione non tentabit Analysta, quo saxa et scrupulos effugiat. ne ignarus \struck{\ill{}} \uncertain{imitetur}, deprimet, attollet et undiq; operabitur semper, nova quæ \ill{} postulaverit \uncertain{suscepto} Zetesi sub alio quam qui proponebatur termino, ad propositam tamen habenti datam differentiam vel rationem.
\note{Le mot barré après « ignarus » commence par « Anat… ». À droite de « rationem. », un signe à l'encre noire et épaisse, en forme de crochet, de la même encre que le « 1 » du haut de la page.}

Omnino Æquationum \uncertain{ortus} et prima constitutio scitu digna est, \uncertain{illaq;} non solita ab Analysta repetenda et adsignanda quo sibi pateat ad eas \uncertain{deductionis} via.

Æqualitatum constitutio potissimo deprehenditur, Zetesi, plasmate, et Syncrisi.
\note{Le texte s'arrête au bas du premier tiers inférieur ; le reste de la page est blanc.}

\page{6}
\note{Verso du feuillet de la vue 5 : en haut à droite, à l'envers et en transparence, « Libri 1201 » et son accolade (lus sur la découpe retournée) ; en haut à gauche, la tache noire du « 1 » du recto, et en bas à droite l'accolade de « R.c. 8070 (73) », vus de même par transparence. Aucun nombre n'est écrit sur cette page. Même main et même encre pâle qu'à la vue 5. Le texte s'arrête un peu au-dessous du milieu de la page.}

\subsection*{De Zetesi Caput II}

Zetesin non instituit \uncertain{tendere} neq; \struck{at} ἀτέχνως analysta, tam \uncertain{pura} \uncertain{\& paris} affecta \& \uncertain{affectis} \uncertain{prosequi} \uncertain{dictat} ratio. Et ideo ad deprehen\add{de}ndam æquationum ad potestates puras pertinentium conditionem, is \uncertain{qui} datis duobus lateribus potestates in genere, ad æquationes affectarum simpliciter potestatum Investigabit \& data differentia vel summa laterum, vel suorum graduum cum dato sub iis \ill{} eorum gradibus paribus vel disparibus facto, alterutrum è lateribus. Ad æquationes denique multipliciter affectarum, affectionem affectioni adiungens, sub dimidiis laterum parodicis gradibus.
\note{La syllabe « de » de « deprehendendam » est ajoutée au-dessus de la ligne par l'auteur, avec un signe d'insertion. Dans la première ligne, un long trait horizontal part de « non » et traverse « instituit » et le mot suivant : c'est d'abord la barre des t, et il n'est pas transcrit comme une rature ; seul « at », repassé d'un trait plus épais, est barré. Le mot grec est en caractères grecs dans la ligne. La deuxième ligne, lue « tam pura \& paris affecta \& affectis prosequi dictat ratio », est d'une lecture très incertaine.}

\uncertain{At} etiam non se adiget uni quæstionis formulæ, sed se sub Zeteticis exercebit sub quocunq; effecto vel proposito themate.

Quæ Geometrico operi \uncertain{aptamus}, recordabimur latera ad extremas lineas rectas in serie continue proportionalium referri, facta vero sub lateribus, aut laterum paribus vel disparibus gradibus, ad aliquam è mediis potestates.

Et cum incidet in æquationem \uncertain{mechanico} sive bene obviam, de illa condet theorema simile \struck{\ill{}} cuiuscumq; æqualitatis Systaticum et Ecdoticum.
\note{Le mot barré après « simile » est couvert d'un trait épais.}

Æquationem Æquationi similem regulariter dicimus cum utrobiq; par est potestas, seu æq. alta, ipsaq; affecta affectionibus sub pari gradu, et eadem affectionis nota. Quod si aliud quiddam præterea requiritur speciale et conditionarium, similitudo anomala est.

\subsection*{Constitutiva Æquationum quadraticarum è Zeteticis Cap. III.}

Affectarum sane quæ \uncertain{adæquant} quadraticarum constitutio è Zeteticis dignoscitur probè, sunt autem Æquationum huiusmodi tres species καταφατική, ἀποφατική, ἀμφίβολος.

Tribus itaq; theorematis de earum constitutione, Analysta ita poterit ratiocinari.
\note{Les trois noms d'espèces sont écrits en grec dans la ligne. La page s'arrête ici ; la suite est à la vue 7.}

\page{7}
\note{Recto d'un nouveau feuillet, même main et même encre pâle que les vues 5–6. En haut à droite, à l'encre, le chiffre « 2 ». Au bord gauche de la vue paraît, sous ce feuillet, la marge d'un autre feuillet portant des fins de mots, qui ne sont pas transcrites. Le texte occupe les deux tiers supérieurs de la page ; le bas est blanc, le bord inférieur abîmé.}

\subsection*{Καταφατικῆς}

\subsection*{Theorema 1.}

Si $A$ quadratum plus $B$ in $A$ æquetur $Z$ quadrato; sunt tres proportionales radices, quarum media est $Z$, differentia vero extremarum $B$. Et fit $A$ minor extrema.

Ex Zetetico \uncertain{si placet} data media trium proportionalium linearum rectarum, et differentia inter extremas, Invenire minorem extremam.

quæ enim in lineis rectis locum habent comparationes easdem ad radices easdem quascumq; simplices, planas, solidasve, aut ulterioris ordinis homogeneas posse aptari \uncertain{disserunt} Campanus in libro de proportione.

Sit igitur data $Z$ media trium proportionalium, $B$ differentia extremarum, et oporteat Invenire minorem extremam. Esto illa $A$. maior igitur erit $A + B$. ducatur minor in maiorem fit $A$ quadratum plus $B$ in $A$. Et vero quando sunt proportionales, quod fit sub extremis æquatur mediæ quadrato: igitur $A$ quadratum plus $B$ in $A$, æquale est $Z$ quadrato. id autem ipsum est quod ordinatur.
\note{« sunt proportiona\-les » est repassé à l'encre plus noire en fin de ligne, de même que « minorem » plus haut, écrit sur un premier mot.}

$1Q + 10N.$ æquatur $144$. Est $L.\,144$ media inter extremas, \uncertain{differentia} est $10$. et fit $1N.$ minor extrema seriei trium proportionalium $8.\ 12.\ 18.$
\note{« L. 144 » : le côté (la racine) de 144, soit 12, le moyen terme de la série 8, 12, 18. $Q$ et $N$ sont le carré et le nombre (la racine) de l'exemple numérique ; le $Q$ a la forme d'un 2 souligné.}

\subsection*{Ἀποφατικῆς}

\subsection*{Theorema II}

Si $A$ quadratum minus $B$ in $A$ æquetur $Z$ quadrato: sunt tres proportionales, quarum media est $Z$, differentia vero extremarum est $B$, et fit $A$ maior extrema.

Ex Zetetico, data media trium proportionalium, et differentia inter extremas Invenire maiorem extremam. Esto illa $A$. minor igitur erit $A - B$. ducatur maior in minorem fit $A$ quadratum minus $B$ in $A$ æquale $Z$ quadrato. Id autem ipsum est quod ordinatur.

$1Q - 10N.$ æquatur $144$: est $L.\,144$ media inter extremas \uncertain{differentia} est $10$ et fit $1N$ maior extrema è serie trium proportionalium $8.\ 12.\ 18.$

\page{8}
\note{Verso du feuillet de la vue 7, sans nombre. Même main. Le texte suit celui de la vue 7 sans rupture et remplit la page jusqu'au bord inférieur, abîmé. Dans la marge de gauche, en regard de la deuxième ligne, un petit signe pâle en forme de « 5 » ou de crochet.}

\subsection*{Ἀμφιβόλου.}

\subsection*{Theorema III.}

Si $B$ in $A$ $-$\add{minus} $A$ quadrato, æquetur $Z$ quadrato; sunt tres proportionales, quarum media est $Z$ aggregatum extremarum $B$ et fit $A$ maior minorve extrema.
\note{Au-dessus du signe « $-$ », l'auteur a écrit « minus » en petites lettres.}

Ex Zetetico, data media trium proportionalium et aggregato extremarum Invenire alterutram extremam.

Sit enim data $Z$ media aggregatum vero extremarum $B$ oportet Invenire minorem extremam. esto illa $A$, maior igitur erit $B - A$ quare $B$ in $A$ minus $A$ quadrato, æquabitur $Z$ quadrato. Sed sit $A$ maior extremarum erit $B - A$ minor extremarum Atq; rursus $B$ in $A$, minus $A$ quadrato æquabitur $Z$ quadrato. unde $A$ sive de minore extremarum, sive de maiore potest enunciari.

$26N. - 1Q$ æquatur $144$. est $L.\,144$ media inter extremas, quarum aggregatum est $26$, et fit $1N.$ minor, maiorve extrema è serie trium proportionalium. $8.\ 12.\ 18.$

\subsection*{Rursus Ex Zeteticis}

\subsection*{Constitutiva æquationum Cubicarum ac primum earum in quibus affectiones desinunt sub latere. Cap. IIII.}

Æquationum quoq; cubicarum affectionibus sub latere involutarum constitutio è Zeteticis scitu digna est, quo pertinent tria quæ sequuntur Theoremata.

\subsection*{Καταφατικῆς}

\subsection*{Theorema I.}

Si $A$ cubus plus $B$ quadrato in $A$, æquetur $B$ quadrato in $Z$. sunt quatuor continue proportionales quarum prima, maior minorve inter extremas est $B$, aggregatum vero secundæ et quartæ est $Z$. et fit $A$ secunda.

Ex Zetetico data prima et aggregato secundæ et quartæ in serie quatuor continue proportionalium, Invenire secundam.

Esto secunda $A$ quarta igitur erit $Z - A$. solido autem sub primæ quadrato et quarta æquatur cubus è secunda, cum sit ut quadratum primæ, ad quadratum secundæ, ita secunda ad quartam. itaq; $A$ cubus æquabitur $B$ quadrato in $Z$, minus $B$ quadrato in $A$. et per antithesin $A$ cubus plus $B$ quadrato in $A$ æquabitur $B$ quadrato in $Z$, ut est ordinatum.
\note{« involutarum » : la première syllabe est à demi effacée au début de la ligne. La page finit ici ; la suite est au recto suivant, vue 9.}

\page{9}
\note{Recto d'un nouveau feuillet, même main ; en haut à droite, à l'encre, le chiffre « 3 ». Le texte continue sans rupture l'exemple du Théorème I, énoncé au bas de la vue 8. L'écriture du verso transparaît fortement dans la marge de gauche.}

Si $1C + 64N.$ æquetur $2496$; sunt quatuor continue proportionales quarum prima minor inter extremas est $L\,64$, id est $8$. aggregatum vero secundæ et quartæ $\dfrac{2496}{64}$, id est $39$ et fit $1N$ secunda è serie proportionalium
\[ \overset{\mathrm{I}}{8}.\ \overset{\mathrm{II}}{12}.\ \overset{\mathrm{III}}{18}.\ \overset{\mathrm{IIII}}{27}. \]
\note{Au-dessus des quatre termes, les rangs I, II, III, IIII en chiffres romains, comme ici.}

Et si $1C + 729N$ æquatur $18954$ prima maior inter extremas est $L\,729$, id est $27$, aggregatum secundæ et quartæ $\dfrac{18954}{729}$ id est $26$. et fit \struck{$x$} $1N$ secunda è eadem serie.

\subsection*{Ἀποφατικῆς}

\subsection*{Theorema II}

Si $A$ cubus minus $B$ quadrato in $A$, æquetur $B$ quadrato in $D$, sunt quatuor continue proportionales, quarum prima minor inter extremas est $B$. differentia secundæ et quartæ $D$ et fit $A$ secunda.

Ex Zetetico. data prima minore inter extremas et differentia secundæ et quartæ in serie quatuor continue proportionalium Invenire secundam.

Sit data $B$ prima minor inter extremas, differentia vero secundæ et quartæ $D$. oportet Invenire continue proportionales. sit secunda $A$ quarta igitur erit $A + D$. Solido autem sub quadrato primæ et quarta, æquatur cubus è secunda, quare $A$ cubus, æquabitur $B$ quadrato in $A$ plus $B$ quadrato in $D$. et per Antithesin $A$ cubus minus $B$ quadrato in $A$ æquabitur $B$ quadrato in $D$. ut est ordinatum.
\note{Après « prima », une tache d'encre couvre une lettre, sans doute un $B$ repassé.}

Si $1C - 64N.$ æquetur $960$. sunt quatuor continue proportionales quarum prima est \struck{latus} $L\,64$. id est $8$. differentia vero secundæ et quartæ $\dfrac{960}{64}$ id est $15$. et fit $1N.$ secunda Ex serie proportionalium $\overset{\mathrm{I}}{8}.\ \overset{\mathrm{II}}{12}.\ \overset{\mathrm{III}}{18}.\ \overset{\mathrm{IIII}}{27}.$

\subsection*{Ἀμφιβόλου}

\subsection*{Theorema III.}

Si $B$ quadratum in $A$ minus $A$ cubo æquetur $B$ quadrato in $D$; sunt quatuor continue proportionales, quarum prima maior inter extremas est $B$. differentia vero secundæ et quartæ $D$. et fit $A$ secunda.
\note{Le texte s'arrête ici, au milieu de la page ; le Zététique du Théorème III suit au verso, vue 10.}

\page{10}
\note{Verso du feuillet de la vue 9, sans nombre. Même main. L'écriture du recto transparaît dans la marge de gauche ; ce n'est pas une note. Le texte s'arrête aux trois cinquièmes de la page ; le bas est blanc.}

Ex Zetetico data prima maiore inter extremas, et differentia secundæ et quartæ, in serie quatuor continue proportionalium Invenire secundam.

Sit data $B$ prima maior inter extremas, differentia vero secundæ et quartæ $D$. oportet Invenire secundam. Sit secunda $A$. quarta igitur erit $A - D$. Solido autem sub quadrato primæ et quarta, æquatur cubus è secunda. quare $A$ cubus æquabitur $B$ quadrato in $A$, minus $B$ quadrato in $D$. et per Antithesin $B$ quadratum in $A$, minus $A$ cubo, æquabitur $B$ quadrato in $D$.

Si $729N - 1C.$ æquentur $7290$; sunt quatuor continue proportionales, quarum prima est $L.\,729$. differentia vero secundæ et quartæ $\dfrac{7290}{729}$ id est $10$ et fit $1N$ secunda è serie proportionalium. $\overset{\mathrm{I}}{27}.\ \overset{\mathrm{II}}{18}.\ \overset{\mathrm{III}}{12}.\ \overset{\mathrm{IIII}}{8}.$

potest autem ea secunda duplex esse, ut in duplici quatuor continue proportionalium serie quæ sequitur.
\[ \begin{array}{llll} \overset{\mathrm{I}}{L.\,59319}. & \overset{\mathrm{II}}{195}. & \overset{\mathrm{III}}{L\,24375}. & \overset{\mathrm{IIII}}{125}. \\[4pt] L\,59319. & 78. & L\,624. & 8. \end{array} \]

Stante eadem prima maiore inter extremas $L.\,59319$. et eadem differentia secundæ et quartæ $70$. secunda hic fit $78$. illic $195$.

Sit stante eadem prima $36$, eademq; differentia secundæ et tertiæ $5$. proponitur duplex series trium proportionalium
\[ \begin{array}{lll} \overset{\mathrm{III}}{36}. & \overset{\mathrm{II}}{6}. & \overset{\mathrm{I}}{1}. \\[4pt] 36. & 30. & 25. \end{array} \]
\note{Dans la première série de trois, les rangs sont écrits à rebours, III sur 36, II sur 6, I sur 1.}

\subsection*{Constitutiva æquationum Cubicarum in quibus affectiones sunt sub quadrato. Cap. V.}

Quæ autem cubicæ affectionibus sub quadrato obvolvuntur æquationes, iisdem fere terminis constant, quæ in affectionibus sub latere. ut è Zeteticis similiter clarum fit. quo pertinebunt trina quoq; item enumeranda theoremata.
\note{La page s'arrête sur cette annonce ; les théorèmes suivent au recto suivant, vue 11.}

\page{11}
\note{Recto d'un nouveau feuillet, même main ; en haut à droite, à l'encre, le chiffre « 4 ». Au bord gauche de la vue, la marge droite d'un autre feuillet porte des fins de lignes qui ne sont pas transcrites. Le texte s'arrête aux deux tiers de la page.}

\subsection*{Ἀποφατικῆς}

\subsection*{Theorema I.}

Si $A$ cubus minus $B$ in $A$ quadrato, æquetur $B$ in $Z$ quadratum; sunt quatuor continue proportionales, quarum prima maior, minorve inter extremas est $B$, aggregatum vero secundæ et quartæ est $Z$. et fit $A$ aggregatum primæ et tertiæ.

Ex Zetetico data prima et aggregato secundæ et quartæ in serie quatuor continue proportionalium, Invenire aggregatum primæ et tertiæ.

Sit data prima $B$. maior minorve inter extremas: aggregatum vero secundæ et quartæ $Z$. in serie quatuor continue proportionalium. oportet Invenire aggregatum primæ et tertiæ.

Sit illud $A$, tertia igitur erit $A - B$ est autem ut $A$ ad $Z$ ita $B$ ad $\dfrac{B \text{ in } Z}{A}$ quare $\dfrac{B \text{ in } Z}{A}$ erit secunda cum sit ut aggregatum primæ et tertiæ ad aggregatum secundæ et quartæ ita prima ad secundam. at rectangulum sub prima et tertia æquabitur secundæ quadrato. Itaq; $\begin{array}{l} B \text{ in } A \\ - B \text{ quadrato} \end{array}$ æquabitur $\dfrac{B \text{ quadrato in } Z \text{ quadratum}}{A \text{ quadratum}}$ omnia multiplicentur per $A$ quadratum et dividantur per $B$. igitur $\begin{array}{l} A \text{ cubus} \\ - B \text{ in } A \text{ quadrato} \end{array}$ æquabitur $B$ in $Z$ quadrato. id autem est quod ordinatum.
\note{Les expressions en deux lignes sont écrites ainsi, l'une au-dessus de l'autre ; la seconde est fermée à droite par une accolade. Devant la fraction « $B$ quadrato in $Z$ quadratum », une tache d'encre.}

Si $1C - 8Q$ æquetur $12168$. sunt quatuor continue proportionales quarum prima inter extremas est $8$ aggregatum vero secundæ et quartæ $L\,\dfrac{12168}{8}$ id est $39$. et fit $1N.$ aggregatum primæ et tertiæ è serie proportionalium
\[ \overset{\mathrm{I}}{27}.\ \overset{\mathrm{II}}{18}.\ \overset{\mathrm{III}}{12}.\ \overset{\mathrm{IIII}}{8}. \]
\note{Le $Q$ de « $8Q$ » a la forme d'un 2 souligné, comme à la vue 7. La série est écrite en commençant par 27, le rang I sur 27, bien que la prima de l'exemple soit 8.}

\subsection*{Καταφατικῆς}

\subsection*{Theorema II}

Si $A$ cubus plus $B$ in $A$ quadratum æquetur $B$ in $D$ quadrato sunt quatuor continue proportionales, quarum prima minor inter extremas est $B$. differentia vero secundæ et quartæ est $D$. et fit $A$ differentia primæ et tertiæ.

\page{12}
\note{Verso du feuillet de la vue 11, sans nombre. Même main. L'écriture du recto transparaît dans la marge de gauche. Dans cette marge, en regard de la ligne « æquabitur $B$ in $Z$ quadratum », une lettre « D » capitale à l'encre noire et épaisse, de la même encre que les chiffres des rectos et que le signe noir de la vue 5 ; sa fonction n'est pas dite. La page est pleine jusqu'au quart inférieur.}

Ex Zetetico data prima minore inter extremas et differentia secundæ et quartæ in serie quatuor continue proportionalium Invenire differentiam primæ et tertiæ.

Sit data $B$ prima in serie quatuor continue proportionalium eademq; minor inter extremas. differentia vero secundæ et quartæ $D$. oportet Invenire differentiam primæ et tertiæ. Esto illa $A$. tertia igitur erit $A + B$. est autem ut $A$ ad $D$ ita $B$ ad $\dfrac{B \text{ in } D}{A}$ quare $\dfrac{B \text{ in } D}{A}$ erit secunda. cum sit ut differentia primæ et tertiæ ad differentiam secundæ et quartæ ita prima ad secundam. at rectangulum sub prima et tertia æquatur secundæ quadrato, Itaq; $\left.\begin{array}{l} B \text{ in } A \\ + B \text{ quadrato} \end{array}\right\}$ æquabitur $\dfrac{B \text{ quadrato in } D \text{ quadrato}}{A \text{ quadrato}}$ omnia ducantur in $A$ quadratum et dividantur per $B$. igitur $A$ cubus, plus $B$ in $A$ quadratum æquabitur $B$ in $Z$ quadratum. id autem ipsum est quod ordinatum.
\note{« $B$ in $Z$ quadratum » : le $Z$ est souligné sur le feuillet ; l'énoncé du Théorème II (vue 11) et le calcul qui précède portent $D$.}

Si $1C + 8Q$ æquetur $1800$. sunt quatuor continue proportionales, quarum prima inter extremas est $8$. differentia vero secundæ et quartæ $L.\,\dfrac{1800}{8}$ id est $15$ et fit $1N.$ differentia primæ et tertiæ è serie proportionalium.
\[ \overset{\mathrm{I}}{8}.\ \overset{\mathrm{II}}{12}.\ \overset{\mathrm{III}}{18}.\ \overset{\mathrm{IIII}}{27}. \]

\subsection*{Ἀμφιβόλου.}

\subsection*{Theorema III.}

Si $B$ in $A$ quadratum minus $A$ cubo æquetur $B$ in $D$ quadrato sunt quatuor continue proportionales quarum prima maior inter extremas est $B$. differentia vero secundæ et quartæ est $D$. et fit $A$ differentia primæ et tertiæ.

Ex Zetetico data prima maiore inter extremas et differentia secundæ et quartæ in serie quatuor continue proportionalium Invenire differentiam secundæ et tertiæ.
\note{« secundæ et tertiæ » est bien la leçon du feuillet ; l'énoncé qui précède et la phrase qui suit portent « primæ et tertiæ ».}

Sit data $B$ prima in serie quatuor continue proportionalium eademq; maior inter extremas, differentia vero secundæ et quartæ $D$. oportet Invenire differentiam primæ et tertiæ. Esto illa $A$ tertia igitur erit $B - A$. Erit autem ut $A$ ad $D$ ita $B$ ad $\dfrac{B \text{ in } D}{A}$. quare $\dfrac{B \text{ in } D}{A}$ erit secunda. cum sit ut differentia primæ et tertiæ ad differentiam secundæ et quartæ ita prima ad secundam at rectangulum sub prima et tertia æquatur secundæ quadrato. Itaq; $\left.\begin{array}{l} B \text{ quadrato} \\ - B \text{ in } A \end{array}\right\}$ æquabitur $\dfrac{B \text{ quadrato in } D \text{ quadratum}}{A \text{ quadrato}}$
\note{La page s'arrête sur cette égalité ; la suite du calcul est attendue au recto suivant. Le « t » initial de « tertiæ », au début de l'avant-dernière ligne, est repassé à l'encre noire.}

\page{13}
\note{Recto d'un nouveau feuillet, même main pour le texte ; en haut à droite, à l'encre, le chiffre « 5 ». Le texte continue sans rupture le Théorème III arrêté au bas de la vue 12. Au pied de la page, sous un trait ondulé tiré au travers, une addition d'une autre main, plus petite et plus rapide, à l'encre plus noire, avec en marge « vide Zet. 20 lib. 2. » ; elle est transcrite à la fin de la vue.}

omnia ducantur in $A$ quadratum et per $B$ dividantur. Igitur $\left.\begin{array}{l} B \text{ in } A \text{ quadratum} \\ - A \text{ cubo} \end{array}\right\}$ æquabitur $B$ in $D$ quadrato. Id autem ipsum est quod ordinatum.
\note{Le $D$ de « $B$ in $D$ quadrato » est repassé en une tache d'encre noire.}

Si $27Q - 1C$ æquentur $2700$. sunt quatuor continue proportionales, quarum prima maior inter extremas est $27$. differentia vero secundæ et quartæ $L\,\dfrac{2700}{27}$ id est $10$ et fit $1N.$ differentia primæ et tertiæ è serie proportionalium
\[ \overset{\mathrm{I}}{27}.\ \overset{\mathrm{II}}{18}.\ \overset{\mathrm{III}}{12}.\ \overset{\mathrm{IIII}}{8}. \]

\subsection*{Alia insuper Constitutio Cuborum sub latere adfectorum. Ex Zeteticis Cap. VI.}
\note{Le « A » de « Alia » est repassé à l'encre noire.}

Sed neq; omittenda est è Zeteticis quoq; depromenda singularis quædam constitutio. cubi, qui sub latere affirmate afficitur, atq; etiam eius qui adficitur negate, cum videlicet (is enim casus præmonendus est) quadruplus cubus è triente coefficientis plani \uncertain{cedit} solidi dati \uncertain{medietatis} quadrato. quocirca bina iam proferuntur theoremata.

\subsection*{Theorema I.}

Si $A$ cubus plus $B$ plano ter in $A$, æquetur $D$ solido, est $B$ planum quod fit sub lateribus, a quibus qui fiunt cubi differunt per $D$ solidum. et fit $A$ differentia laterum.

Ex Zetetico, data differentia cuborum et rectangulo sub lateribus Invenire differentiam laterum.

Si $1C + 6N$ æquatur $7$. est $\dfrac{6}{3}$ seu $2$ rectangulum sub lateribus a quibus cubi differunt per $7$. et fit $1N$ differentia laterum. è hypothesi laterum $1$ et $2$.

\subsection*{Theorema II.}

Si $A$ cubus minus $B$ plano ter in $A$ æquetur $D$ solido. præstet autem $D$ solidi quadratum quadruplo $B$ plani cubo. est $B$ planum rectangulum sub lateribus a quibus qui fiunt cubi componunt $D$ solidum et fit $A$ aggregatum laterum.

Ex Zetetico. dato rectangulo sub lateribus et aggregato cuborum Invenire latera.
\note{« et aggregato » est d'une encre plus noire que la ligne, peut-être repassé ou ajouté.}

\marginal{vide Zet. 20 lib. 2.}
adscripta conditio hic patet. sit latera $B$. $A$. rectangulum sub lateribus $B$ in $A$. aggregatum cuborum $A\,c + B\,c$. cubus rectanguli sub lateribus est $B\,c$ in $A\,c$. quadratum aggregati cuborum $\begin{array}{l} A\,cc \\ + B\,c \text{ in } A\,c\ 2 \\ + B\,cc. \end{array}$ est autem $\begin{array}{l} A\,cc \\ + B\,c \text{ in } A\,c\ 2 \\ + B\,cc \end{array}$ maius $B\,c$ in $A\,c\ 4$. quandoquidem $A\,cc + B\,c$ in $A\,c. + B\,cc.$ sunt continue proportionales atq; in tribus continue proportionalibus in ratione inæqualitatis summa extremarum maior est mediæ dupla. quare $A\,cc + B\,cc$ maius erit $B\,c$ in $A\,c\ 2$. unde constat propositum. \uncertain{quando} \ill{}
\note{Cette addition est d'une autre main que le texte, plus petite, à l'encre noire ; elle commence plus à gauche que la justification et déborde la marge. La note marginale « vide Zet. 20 lib. 2. » est de la même main. La première expression à trois lignes est entourée d'un trait qui la rattache à la fin de la deuxième ligne. Dans cette main, « $A\,c$ », « $A\,cc$ » s'écrivent avec des $c$ qui ressemblent à des $r$. Le dernier mot, après « quando », est sous une tache au bord déchiré du feuillet.}

\page{14}
\note{Verso du feuillet de la vue 13, sans nombre. Au-dessus de son bord supérieur paraissent trois lignes d'un autre feuillet : ce sont les premières lignes de la vue 12 (« Ex Zetetico data prima minore inter extremas … Invenire differentiam primæ et tertiæ. »), que le feuillet 5, plus court, laisse voir ; elles ne sont pas transcrites ici. Le texte continue le Théorème II de la vue 13. Deux notes marginales à gauche sont de la main de l'addition de la vue 13, à l'encre noire.}

\marginal{$B\,p.$ rectangulo sub lateribus / $D\,s.$ aggregato cuborum / quadruplus $B\,p$ cubus. \uncertain{cedit} $D$ solidi quadrato}
Si $1C - 6N.$ æquatur $9$ est $\dfrac{6}{3}$ seu $2$. rectangulum sub lateribus a quibus cubi aggregati faciunt $9$ et fit $1N$ aggregatum laterum è hypothesi laterum $1$ et $2$.
\note{La note marginale, en quatre lignes, est reliée par une accolade au « Si » qui ouvre l'exemple ; elle nomme les données du Théorème II et répète sa condition.}

Eadem autem theoremata Geometrice ita concipiuntur.

\subsection*{Aliter. Primum theorema.}

Si $A$ cubus plus $B$ plano ter in $A$ æquetur $B$ plano in $D$ sunt quatuor continue proportionales lineæ rectæ sub quarum mediis vel extremis fit $B$ planum. differentia vero extremarum est $D$. et fit $A$ differentia mediarum.

\struck{Si $1C + 24N.$ æquatur $56$, sunt quatuor continue proportionales}

Ex Zetetico data differentia extremarum et rectangulo sub mediis vel extremis in serie quatuor proportionalium Invenire differentiam mediarum.

Si $1C + 24N.$ æquatur $56$ sunt quatuor continue proportionales sub quarum mediis vel extremis quod fit planum æquale est $\dfrac{24}{3}$ id est $8$ differentia vero extremarum est $\dfrac{56}{8}$ id est $7$ et fit $1N.$ differentia mediarum è serie continue proportionalium $\overset{\mathrm{I}}{1}.\ \overset{\mathrm{II}}{2}.\ \overset{\mathrm{III}}{4}.\ \overset{\mathrm{IIII}}{8}.$
\note{La phrase barrée est un premier départ de l'exemple, écrit avant le Zététique puis rayé d'un trait continu ; l'exemple est récrit après.}

\subsection*{Aliter. Secundum theorema.}

\marginal{est enim aggregatum extremarum in serie quatuor continue proportionalium maius aggregato mediarum quadratum autem semissis mediarum maius est rectangulo sub iisdem sive quadrato mediæ proportionalis inter easdem medias quare et extremarum semissis quadratum eodem rectangulo erit multo maius.}
Si $A$ cubus minus $B$ plano ter in $A$ æquetur $B$ plano in $D$, sit autem $D$ semissis quadratum maius $B$ plano, sunt quatuor continue proportionales lineæ rectæ sub quarum mediis vel extremis fit $B$ planum. aggregatum vero extremarum est $D$. et fit $A$ aggregatum mediarum.

Ex Zetetico dato aggregato extremarum, et rectangulo sub mediis vel extremis in serie quatuor continue proportionalium, Invenire aggregatum mediarum.

Si $1C - 24N.$ æquatur $72$ sunt quatuor continue proportionales, sub quarum extremis vel mediis quod fit planum æquale est $\dfrac{24}{3}$ id est $8$ aggregatum vero extremarum est $\dfrac{72}{8}$ id est $9$ et fit $1N$ aggregatum mediarum è serie continue proportionalium $\overset{\mathrm{I}}{1}.\ \overset{\mathrm{II}}{2}.\ \overset{\mathrm{III}}{4}.\ 8.$
\note{Ici, seuls les trois premiers termes portent un rang.}

Cum autem quadruplus cubus è triente coefficientis plani præstat solidi dati medietatis quadrato, aliam eo casu sortitur æqualitas constitutionem ambiguæ seu eius quæ negatur \uncertain{Inversæ}, communem, quo pertinet sequens theorema.
\note{Le mot lu « Inversæ » montre « Inu…st… » ; la lecture reste douteuse. La page s'arrête sur cette annonce ; le théorème annoncé est attendu au recto suivant.}

\page{15}
\note{Petit feuillet inséré, écrit seulement dans sa moitié supérieure, de la main rapide à l'encre noire qui a fait l'addition de la vue 13 et les notes marginales de la vue 14 — non de la main du texte. En haut à droite, à l'encre, « 5 bis » : le feuillet se donne ainsi pour une insertion après le feuillet paginé « 5 » (vues 13–14). En haut, au milieu, un signe « $\overset{\circ}{x}$ 4 » dont le sens n'est pas établi. Cachet rond de la Bibliothèque au milieu, qui couvre quelques signes. C'est la démonstration de la condition énoncée à la vue 13 (« quadruplus cubus è triente coefficientis plani … »), renvoyée là par « adscripta conditio hic patet ». Dans cette main les puissances s'écrivent $A\,q$, $A\,c$, $A\,qq$, $A\,qc$, $A\,cc$, avec des $c$ qui ressemblent à des $r$ ; un chiffre écrit après une expression en est le coefficient.}

vel sic \uncertain{ἀποδεικτικῶς}

quoniam ex $A + B$ in $B$ in $A$ fit $\begin{array}{l} B\,q \text{ in } A \\ + B \text{ in } A\,q \end{array}$. cuius quadratum si comparetur \struck{cum} \struck{sub} cum quadruplo cubo ab $\dfrac{\begin{array}{l} A\,q \\ + B\,q \\ + B \text{ in } A \end{array}}{3}$ fiunt illinc $\begin{array}{l} B\,qq \text{ in } A\,q\ 3 \\ + B\,q \text{ in } A\,qq\ 3 \\ + B\,c \text{ in } A\,c\ 26 \end{array}$ hinc vero $\begin{array}{l} A\,cc\ 4 \\ + B \text{ in } A\,qc\ 12 \\ + B\,cc\ 4 \\ + B\,qc \text{ in } A\ 12 \end{array}$ (reiectis scilicet è facti solido-solidis \struck{\ill{}} utrinq; æqualibus.) Sunt autem proportionales continue $\begin{array}{l} A\,cc \\ + B \text{ in } A\,qc \\ + B\,q \text{ in } A\,qq \\ + B\,c \text{ in } A\,c \\ + B\,qq \text{ in } A\,q \\ + B\,qc \text{ in } A \\ + B\,cc \end{array}$ quare erit quoq; ut $B$ in $A\,qc$ ad $B\,q$ in $A\,qq$ ita $B\,qq$ in $A\,q$ ad $B\,qc$ in $A$ et summa extremarum præstabit summæ mediarum. quare et $\begin{array}{l} B \text{ in } A\,qc\ 12 \\ + B\,qc \text{ in } A\ 12 \end{array}$ maius erit quam $\begin{array}{l} B\,qq \text{ in } A\,q\ 12 \\ + B\,q \text{ in } A\,qq\ 12 \end{array}$. et ablatis utrinq; $\begin{array}{l} B\,qq \text{ in } A\,q\ 3 \\ + B\,q \text{ in } A\,qq\ 3 \end{array}$. restabit illinc magnitudo maior quam $\begin{array}{l} B\,qq \text{ in } A\,q\ 9 \\ + B\,q \text{ in } A\,qq\ 9 \end{array}$. atqui inter $B\,qq$ in $A\,q$ et $B\,q$ in $A\,qq$ media proportionalis est $B\,c$ in $A\,c$. in serie autem continue proportionalium summa extremarum maior est media dupla. quare $\begin{array}{l} B\,qq \text{ in } A\,q\ 9 \\ + B\,q \text{ in } A\,qq\ 9 \end{array}$ maius erit quam $B\,c$ in $A\,c\ 18$. est autem et $B\,c$ in $A\,c$ media quoq; proportionalis inter $A\,cc$ et $B\,cc$ quare $\begin{array}{l} A\,cc\ 4 \\ + B\,cc\ 4 \end{array}$ maius quoq; erit quam $B\,c$ in $A\,c\ 8$. Summa igitur $\begin{array}{l} A\,cc\ 4 \\ + B \text{ in } A\,qc\ 12 \\ + B\,cc\ 4 \\ + B\,qc \text{ in } A\ 12 \end{array}$ maior erit quam summa $\begin{array}{l} B\,qq \text{ in } A\,q\ 3 \\ + B\,q \text{ in } A\,qq\ 3 \\ + B\,c \text{ in } A\,c\ 26. \end{array}$ quod erat demonstrandum.

unde patet è \uncertain{dat.} hypothesi \ill{} $D$ solidi quadratum \uncertain{cedere} quadruplo $B\,p.$ cubo.
\note{Toutes les sommes sont écrites en colonnes de termes, l'une sous l'autre, comme ici. Le premier mot grec est lu avec doute (« ἀπ…δεικτικῶς »). Après « hypothesi », une tache pâle couvre un ou deux signes. Le reste du feuillet, sous cette ligne, est blanc.}

\page{16}
\note{Recto d'un nouveau feuillet, main du texte ; en haut à droite, à l'encre, le chiffre « 6 ». Le Théorème III annoncé au bas de la vue 14 ouvre la page. La marge de gauche, sur toute la hauteur, et le pied de la page sont couverts d'additions de la main rapide à l'encre noire (celle des vues 13–15), en lignes serrées ; la moitié inférieure de cette colonne est barrée de traits obliques croisés. Cachet rond pâle à droite, sous les premières lignes. Dans le coin inférieur droit, une lettre « D » entourée d'un cercle, de l'encre noire épaisse du « D » marginal de la vue 12.}

\subsection*{Theorema III}

\marginal{Si sit aggregat. laterum $A$ \struck{\ill{}} cuius quadrat. erit $A\,q$. et \struck{$B\,pl$} $A\,q - B$ planum \add{$D$} erit rectangul. sub lateribus æquale $\dfrac{D\,s.}{A}$ quare $\begin{array}{l} A\,c \\ - B\,pl. \text{ in } A \end{array}$ æq. $D\,s.$ / Sit rursus latus unum $A$ cuius \uncertain{ergo} quadrat. erit $A\,q$ et $B\,p - A\,q$ erit summa rectanguli sub lateribus et quadrati reliqui atqui solidum sub aggregato laterum in rectang. sub lateribus applicat. alteri lateri dat rectang. sub aggregato laterum in latus alterum vel aggregat. quadrati lateris alterius et rectanguli sub lateribus \uncertain{eodem} quare $B\,p\,3 - A\,q\,3$ \ill{} $\dfrac{D\,s.}{A}$ ac proinde $\begin{array}{l} B\,p \text{ in } A\,3 \\ - A\,c \end{array}$ \uncertain{æq.} $D\,s.$ et fit $A$ latus alterutrum.}
Si $A$ cubus minus $B$ plano ter in $A$ æquetur $D$ solido, cedat autem $D$ solidi quadratum quadruplo $B$ plani cubo: $B$ planum ter in $E$ minus $E$ cubo æquabitur rursus $D$ solido. et Enunciatur $E$ de duobus lateribus a quibus \uncertain{singulis} quadrata adiecta rectangulo sub ipsis lateribus faciunt $B$ planum. \struck{Ex} $\overset{\circ}{x}\,4$ quod autem fit ab uno laterum in quadratum reliqui, adiunctum rectangulo, seu aliter quod fit abs aggregato laterum in rectangulum est $D$ solidum. $A$ vero enunciatur de aggregato illorum laterum.
\note{La première note marginale, en deux paragraphes séparés par un trait, justifie les deux énoncés ($A$ agrégat des côtés ; $A$ l'un des côtés). Le signe « $\overset{\circ}{x}\,4$ » dans la ligne, après « $B$ planum », est un signe de renvoi ; le même signe se retrouve dans la marge, plus bas, en tête de la seconde addition, et en tête du feuillet « 5 bis » (vue 15), qui commence par « vel sic ».}

Ex Zetetico. dato plano quod constat aggregato quadratorum a duobus lateribus, plus rectangulo sub iisdem, et dato insuper solido, quod fit abs aggregato laterum in rectangulum, Invenire latera vel etiam laterum aggregatum.

Si $1C - 21N.$ æquetur $20$. quoniam quadruplus cubus è $7$ maior est quam $400$. igitur $21N$ $-$ \struck{$1Q$} \add{$1C$} æquabitur $20$. et sunt duo latera, a quibus quadrata adiuncta rectangulo a lateribus, faciunt $21$. aggregatum autem laterum, ductum in rectangulum est $20$. et fit $1N.$ in æqualitate Inversa negata, alterutrum è lateribus, maius minusve, in directa vero negata aggregatum ipsorum laterum. è hypothesi laterum $1.$ et $4.$
\note{« $1C$ » est écrit au-dessus d'un « $1Q$ » barré.}

\marginal{\struck{$\overset{\circ}{x}\,4$ positis duobus lateribus $B$ et $A$. fiet cubus è triente coefficientis plani $\dfrac{\begin{array}{l} A\,cc \\ + A\,qq \text{ in } B\,q\ 6 \\ + B \text{ in } A\,qc\ 3 \\ + B\,qq \text{ in } A\,q\ 6 \\ + B\,c \text{ in } A\,c\ 7 \\ + B\,cc \\ + B\,qc \text{ in } A\ 3 \end{array}}{27}$ at $D\,s.$ quadrat. $\begin{array}{l} B\,qq \text{ in } A\,q \\ + B\,q \text{ in } A\,qq \\ + B\,c \text{ in } A\,c\ 2 \end{array}$ et quadruplato priore reductisq; omnibus in $27$ \uncertain{partes} ablatis communibus remanebit $\begin{array}{l} A\,cc\ 4 \\ + B \text{ in } A\,qc\ 12 \\ + B\,cc\ 4 \\ + B\,qc \text{ in } A\ 12 \end{array}$ quod maius esse debet quam $\begin{array}{l} B\,qq \text{ in } A\,q\ 3 \\ + B\,q \text{ in } A\,qq\ 3 \\ + B\,c \text{ in } A\,c\ 26 \end{array}$ quod vero sic \ill{} patebit quoniam $A\,cc$. $B$ in $A\,qc$. $B\,q$ in $A\,qq$. $B\,c$ in $A\,c$. $B\,qq$ in $A\,q$. $B\,qc$ in $A$. sunt proportionales continue. est autem in serie quotcunq; continue proportionalium summa extremarum maior summa mediarum quare erit primo $B\,qq$ in $A\,q + B\,q$ in $A\,qq$ maior quam $B\,c$ in $A\,c$ \ill{} $B\,qq$ in $A\,q\ 12 + B\,q$ in $A\,qq\ 12$ maior \ill{} quam $B\,c$ in $A\,c\ 12$ atq; $B$ in $A\,qc\ 12 + B\,qc$ in $A\ 12$ maiores \ill{} quam $B\,qq$ in $A\,q\ 12 + B\,q$ in $A\,qq\ 12$ id est multo maiores quam $B\,c$ in $A\,c\ 24$. et $A\,cc\ 4 + B\,cc\ 4$ maior est \add{multo} quam $B\,qq$ in $A\,q\ 4 + B\,q$ in $A\,qq\ 4 + B\,c$ in $A\,c\ 4$ \ill{}}}
\note{Seconde addition marginale, de la même main, qui commence au signe de renvoi « $\overset{\circ}{x}\,4$ » dans la marge et se poursuit au pied de la page sous le texte. Elle est barrée : sa moitié inférieure par de grands traits obliques croisés, ses deux dernières lignes au pied de la page par des traits horizontaux. C'est une première démonstration de la condition « quadruplus cubus è triente coefficientis plani … », que la même main a récrite au propre sur le feuillet « 5 bis » (vue 15, « vel sic ἀποδεικτικῶς »). Dans le produit à sept termes, le « 26 » est écrit sur un premier chiffre ; les deux « 3 » qui précèdent sont écrits sur des « 21 » barrés. La fin, au pied de la page, porte encore, lisibles sous les ratures : « summa extremarum … », « quadruplo », « multo igitur maius erit quadruplum cubi … ». Lecture d'ensemble incertaine.}

Neq; vero in Geometrica phrasi, hic erit magna dissimilitudo. \ill{} dicet Geometra $B$ planum esse aggregatum quadratorum a tribus proportionalibus lineis rectis $D$ vero solidum quod fit ab aggregato extremarum in mediæ quadratum, seu ab alterutra extrema in aggregatum quadratorum è reliquis. et fieri $A$ aggregatum extremarum, $E$ vero primam vel tertiam.
\note{Le mot avant « dicet » montre « ſ…mmendo » ; il n'est pas lu.}

Sic in Exposito themate, $21$ est aggregatum quadratorum a tribus proportionalibus, et solidum $20$ est ab aggregato extremarum in mediæ quadratum, seu ab alterutra extrema in quadrat\struck{um}\add{a} è reliquis. Ex serie proportionalium $1.\ 2.\ 4.$

At Elegantius et præstantius ex analyticis angularium sectionum huiusmodi æqualitatum constitutio eruitur et ad \uncertain{εὐμηχανίαν} accommodatius in hanc formulam.
\note{Le mot grec est en caractères grecs dans la ligne. La phrase annonce une formule qui n'est pas sur cette page ; elle est attendue au verso, vue 17.}

\page{17}
\note{Verso du feuillet de la vue 16, sans nombre. Le texte de la main principale occupe la partie droite ; toute la colonne de gauche et le pied de la page sont de la main rapide à l'encre noire, avec une figure. En tête de cette colonne, dans un cadre, la fin de l'addition barrée du pied de la vue 16, reprise sous le même signe « D » entouré d'un cercle, qui sert de renvoi d'une page à l'autre. La lecture de la colonne marginale est incertaine en beaucoup d'endroits.}

\marginal{(D) \uncertain{a} $\dfrac{\begin{array}{l} B\,q \\ + A\,q \\ + B \text{ in } A \end{array}}{3}$ quam quadrat. solidi $\begin{array}{l} B\,q \text{ in } A \\ + B \text{ in } A\,q \end{array}$ unde constat oppositum \ill{} theorematis 3 antecedentis.}
\note{Figure dans la marge, sous ce cadre : un cercle de diamètre $AB$, avec sur la circonférence, dans l'ordre, les points $A$, $D$, $E$, $F$, $K$, $G$, $B$, $L$, $M$, $H$ ; un point $C$ hors du cercle, à droite, joint à $G$ et à $B$ ; des cordes en pointillé de $A$ vers $D$, $E$, $F$, $K$, $G$, $B$, de $H$ et $M$ vers le haut, et plusieurs intersections intérieures, dont une marquée $H$ au milieu. Un arc $KG$ est repassé à l'encre.}

\subsection*{Aliter, tertium theorema}

Si $A$ cubus, minus $B$ quadrato ter in $A$ æquetur $B$ quadrato in $D$, sit autem \add{$B$} maior $D$ semisse $\left.\begin{array}{l} B \text{ quadratum ter in } E \\ - E \text{ cubo} \end{array}\right\}$ æquabitur $B$ quadrato in $D$.
\note{L'ajout interlinéaire après « autem » est un $B$ surmonté de deux petits signes, peut-être « $B\,q^{o}$ » ; il est donné sous la forme $B$.}

Et sunt duo triangula rectangula æqualis $B$ hypotenusæ, ita ut angulus acutus subtensus a perpendiculo primi sit triplus \add{ad} angulum acutum subtensum a perpendiculo secundi, basis vero dupla primi est $D$. et fit $A$ dupla basis secundi. $E$ vero basis simpla secundi, contracta \ill{} longitudine lineæ quæ potest quadrato triplum perpendiculi eiusdem.
\note{Le mot entre « contracta » et « longitudine » montre « prohartam » ou « pro partem » ; il n'est pas lu.}

$1C - 300N$ æquetur $432$, vel etiam $300N - 1C.$ æquetur $432$. sunt duo triangula rectangula, quorum hypotenusa communis est \struck{\ill{}} $10$ ita ut angulus acutus primi, à perpendiculo videlicet subtensus, sit triplus ad acutum secundi a suo quoq; perpendiculo subtensum, basis autem primi dupla est $\dfrac{432}{100}$. et $1N$ in æqualitate directe negata est basis dupla secundi, In inversa vero negata est basis simpla secundi plus minusve ea quæ potest quadrato triplum perpendiculum secundi.
\note{Dans « $1C - 300N$ », le $N$ est ajouté au-dessus de la ligne.}

Constituta hypotenusa \struck{secundi} communi $10$ basi secundi trianguli $9$ fit perpendiculum eiusdem secundi $L\,19$. primi vero hypotenusa stante $10$ fit basis $2\tfrac{16}{100}$. itaq; cum in ea hypothesi dicetur $1C - 300N.$ æquari $432$, fiet $1N.$ $18$. vel cum dicetur $300N - 1C$ æquari $432$ fiet unus numerus $9 + L\,57$ vel $9 - L\,57$.

Atq; hæc ad assequendum ex Zeteticis constitutiones adæquatarum, quæ affectæ sunt potestatum exempla \uncertain{nunc} \uncertain{sufficiendo}. Etsi enim de simpliciter affectis theoremata tantum proposita sunt, tamen quemadmodum ad multipliciter affectas possint trahi, vix ignorabitur, quando præsertim deteget earum plasma, de quo Jam succedit dicendi locus. ipsorum enim Zeteticorum, a quibus antedictæ constitutiones depromptæ sunt Jam patuit processus sufficienter in Expositis \ill{}, hoc postremo theoremate excepto, quod par est reiici in peculiarem Zeteticorum ad angulares sectiones pertinentium \ill{}.
\note{Le mot après « Expositis » montre « Ex…sil » ; le dernier mot de la page, « Syl…milam », souligné d'un trait qui revient en crochet, n'est pas lu.}

\marginal{Sit circulus cuius diameter $AB$ \uncertain{cuiusque} circumferentiæ \ill{} $FB$ tertia pars sit $GB$. et quoniam \struck{ita} \add{2} demonstratum a nobis ad sectiones angulares \uncertain{Elementa} $GB$, $GK$, $KF$ sunt æquales erit ut semidiameter ad rectam $AG$ ita $AG$ ad compositam è $AB$, $AK$ et ita $AK$ ad compositam è $AG$ $AF$. in rectis simplicibus sit $AB$ ($B$). $AF$ ($D$). $AG$ bis ($A$) ut $B$ semissis ad $\tfrac{A}{2}$ \ill{} $B$ ad $A$ \struck{duplum} ita $A$ ad $\dfrac{A\,q}{B}$ a quo ablato $B$, relinquitur $\dfrac{A\,q - B\,q}{B}$ æqualis ipsi $AK$ \struck{\ill{}} itemque ut $B$ ad $A$ ita $\dfrac{A\,q - B\,q}{B}$ ad $\dfrac{A\,c - B\,q \text{ in } A\ 2}{B\,q}$ a quo sublato $A$ relinquitur $\dfrac{A\,c - B\,q \text{ in } A\ 3}{B\,q}$. æquale sit ipsi $AF$ bis vel $D$. Jam ipsius $BG$ \add{rectæ} dupla fiat \ill{} $BL$. quæ occurrat circumferentiæ in $L$ puncto et quoniam $BG$ tertia pars est circumferentiæ $BF$ et $GB$ tertia pars totius circuli \ill{} nam quoniam $BL$ dupla est ipsius $BG$ erit in triangulo $GBL$ rectangulo angulus $GBL$ $\tfrac{2}{3}$ unius recti sive $\tfrac{1}{3}$ duorum rectorum. et quadratum \ill{} $GL$ triplum quadrati $GB$. tota circumferentia $BD$ metietur circulum integrum et præterea peripheriam $BF$. metiatur peripheriis $BD$, $DM$, $MF$, erit igitur circumferentia $BM$ residua dupli ipsius $BD$ è integro circulo, dupla ipsius $BDA$ residuæ eiusdem è semicirculo, circumferentia vero $BM$ æqualis est ipsa $FD$ quoniam subtensæ $BD$, $DM$ æquales sunt subtensis $DM$, $MF$. circumferentia igitur $DF$ dupla est circumferentiæ $DA$ totaq; $FA$ eiusdem}
\note{Longue addition de la main à l'encre noire : elle occupe la colonne de gauche depuis la figure jusqu'au bas, puis, sous le texte principal, trois lignes sur toute la largeur de la page, qui s'arrêtent au bord inférieur sur « eiusdem » sans conclusion visible ; la suite n'est pas sur ce feuillet. Elle démontre par la figure (trisection de l'arc $FB$ en $G$, $K$) la constitution du « tertium theorema ». Beaucoup de lettres de la figure et de chiffres y sont repassés ou corrigés ; la lecture en est donnée sous réserve. « bis » au-dessus de « $AF$ » et de « $AG$ » est interlinéaire.}

\page{18}
\note{Recto d'un nouveau feuillet, main du texte ; en haut à droite, à l'encre, le chiffre « 7 ». La colonne de gauche porte, sur toute la hauteur, la suite de l'addition à l'encre noire commencée à la vue 17 (démonstration du « tertium theorema » par la figure du cercle) ; sa lecture est incertaine en beaucoup d'endroits. Le feuillet est plus court que le suivant : sous son bord inférieur déchiré paraissent deux lignes et demie d'un autre feuillet (« … altarum, et adfectarum. Æquationis tamen … tarum spectata proprietas maxime Juvat. »), qui ne sont pas transcrites ici.}

\marginal{$DA$ tripla. Secta igitur $DF$ æqualiter in $E$ et ducta $AE$, et quoniam triangula $BIG$, $ADI$ similia sunt erit $AD$ ipsius $AI$ semissis. in rectis itaq; simplicibus sit $AD$ bis sive $AI$ ($E$). $AB$ ($B$). $AF$ bis ($D$) est è demonstratis a nobis \add{\uncertain{in tractatu de sectionibus angularibus}} ut semidiameter ad $AD$ vel $AB$ ad $AI$ hoc est $B$ ad $E$ ita $E\,\tfrac{1}{2}$ ad $\dfrac{E\,q\,\tfrac{1}{2}}{B}$. differentiam scilicet $AB$. $BE$. quo a $B$ ablato relinquitur $\dfrac{B\,q - E\,q\,\tfrac{1}{2}}{B}$ æquale ipsi $BE$. Deinde ut $B$ ad $E$ ita $\dfrac{B\,q - E\,q\,\tfrac{1}{2}}{B}$ ad $\dfrac{B\,q \text{ in } E - E\,c\,\tfrac{1}{2}}{B\,q}$ differentiam $AD$ \ill{} $AF$ \struck{cui adiectum data} $E\,\tfrac{1}{2}$ \struck{\ill{}} dat $\dfrac{B\,q \text{ in } E\,\tfrac{3}{2} - E\,c\,\tfrac{1}{2}}{B\,q}$ æquale ipsi $AF$. cuius duplum $\dfrac{B\,q \text{ in } E\,3 - E\,c}{B\,q}$ \add{vel $D$} æquale ipsi $AF$ bis. / Protrahatur $AG$ in $C$ et fiat $BC$ dupla ipsius $BG$ erit igitur angulus $GBC$ $\tfrac{2}{3}$ unius recti et quadratum $GC$ triplum quadrati $GB$. et angulus $ABH$ æqualis ipsis $ACB$ tertiæ parti unius recti et $CAB$ tertiæ parti circumferentiæ $BF$. est igitur circumferentia $AH$ tertia pars totius circumferentiæ $A$\uncertain{$H$}$E$. in rectis igitur simplicibus sit $AB$ ($B$) $AC$ ($E$). $AF$ bis ($D$) et quoniam triangula $GBC$, $CAH$ sunt similia erit $AH$ semissis ipsius $AC$. erit igitur ex demonstratis a nobis ut $B$ ad $E$ ita $E\,\tfrac{1}{2}$ \uncertain{ad} $\dfrac{E\,q\,\tfrac{1}{2}}{B}$ quo è diametro ablato relinquitur $\dfrac{B\,q - E\,q\,\tfrac{1}{2}}{B}$ æquale ipsi $BL$. deinde ut $B$ ad $E$ ita $\dfrac{B\,q - E\,q\,\tfrac{1}{2}}{B}$ ad $\dfrac{B\,q \text{ in } E - E\,c\,\tfrac{1}{2}}{B\,q}$ differentiam ipsarum $AH$, $AF$. quo ad $AH$ vel $E\,\tfrac{1}{2}$ addito fit $\dfrac{B\,q \text{ in } E\,\tfrac{3}{2} - E\,c\,\tfrac{1}{2}}{B\,q}$ æquale ipsi $AF$ et omnibus duplicatis $\dfrac{B\,q \text{ in } E\,3 - E\,c}{B\,q}$ æquabitur \add{$D$ vel} $AF$ bis quod erat demonstrandum.}
\note{L'addition marginale est donnée sous réserve : les fractions y sont écrites de façon serrée, avec des « ½ » et des « 3/2 » en exposant à côté des lettres, et plusieurs lettres de la figure y sont reprises. Elle se termine par « quod erat demonstrandum », au niveau de la dernière ligne du texte principal.}

\section*{De Generali methodo transmutandarum æquationum Cap VII.}

Plasmatis ratio pendet è doctrina transmutandarum æquationum generaliter Imprimis proponenda et demonstranda.

\subsection*{De transmutatione per alterationem radicis.}

Æquationis transmutatio instituitur duplici formæ \uncertain{permutatione}. aut \struck{enim} placet alterari radicem de qua primo quæritur aut manere Invariatam.

qualiscumq; autem alteratio sit, primigenia radix et nova datam habent inter se differentiam vel rationem, adeo ut una cognita, altera non possit ignorari.

ad opus itaq; transmutationis alterata radix quæsititia conceditur alteri quoq; inquirendæ esse æqualis, atq; adeo è ea concessione novam induit speciem, sub qua æquatio primo posita dirigitur, et ordinatur nova.

ac pluribus quidem modis radix de qua quæritur potest artificiose alterari, et nova specie exhiberi. at dirigendi ratio sub ea nova specie ipsam quæ primo proposita est æquationem, constans est et uniformis per modos illos \uncertain{quosvis}.

quæ ut fiant evidentiora, proponatur æquatio quævis de $A$ radice exponenda, et sit $Z$ magnitudo cui adæquetur \ill{}, et oporteat propositam illam æquationem alterata radice arte transmutare.

primo igitur $A$ radix potest alterari et nova specie exhiberi per additionem, utpote tali concessione et argumentatione $A + B$ esto $E$. ergo $E - B$ erit $A$.

secundo per subductionem, utpote tali concessione et argumentatione $A - B$ esto $E$. ergo $E + B$ erit $A$. vel etiam ista $B - A$ esto $E$ ergo $B - E$ erit $A$.

tertio per multiplicationem. utpote tali concessione et argumentatione $B$ in $A$ esto $E$ planum. ergo $\dfrac{E \text{ planum}}{B}$ \struck{applicatum} erit $A$.

quarto per divisionem. utpote concedendo et argumentando. $\dfrac{A \text{ planum}}{B}$ esto $E$. ergo $B$ in $E$ erit $A$ planum.

quinto per analogiam rationis explicatæ. utpote concedendo esse ut $B$ ad $G$ ita $A$ ad $E$, deinde resolvendo analogiam et argumentando ergo $\dfrac{B \text{ in } E}{G}$ erit $A$.
\note{Le texte de ce feuillet s'arrête ici, au-dessus du bord inférieur déchiré ; la suite est attendue au verso, vue 19.}

\page{19}
\note{Verso du feuillet de la vue 18, sans nombre ; main du texte seule, sans addition dans la marge de gauche (l'écriture du recto y transparaît). Le texte continue l'énumération des modes d'altération de la racine commencée à la vue 18. Au-delà du bord droit de la page, près de la reliure, sur un talon ou sur le feuillet voisin, trois fragments d'une main plus petite, coupés : « + in ecomp… » (au niveau de la sixième ligne), « vel potius … » et, plus bas, « + vel in … » ; ils ne sont pas transcrits. La page s'arrête aux trois quarts.}

Sexto per analogiam rationis explicitæ. utpote concedendo esse ut $A$ ad $B$ ita $G$ ad $E$, deinde resolvendo analogiam et argumentando ergo $\dfrac{B \text{ in } G}{E}$ erit $A$.

Septimo per parabolicam hypostasin in \ill{} quibuscumq; æquationibus, utpote in quadraticis concedendo $\left.\begin{array}{l} E \text{ quadratum} \\ + A \text{ in } E \end{array}\right\}$ æquari $D$ plano, quæ æquatio est quadrati affirmate affecti deinde argumentando ergo $\dfrac{D \text{ plano} - E \text{ quadrato}}{E}$ erit $A$.
\note{Le mot après « in » montre « condicibus » ou « ordinibus » ; il est repassé.}

vel concedendo $\left.\begin{array}{l} E \text{ quadratum} \\ - A \text{ in } E \end{array}\right\}$ æquari $D$ plano, quæ æquatio est quadrati negate affecti. et argumentando ergo $\dfrac{E \text{ quadrato} - D \text{ plano}}{E}$ æquabitur $A$.

vel denique concedendo $\left.\begin{array}{l} E \text{ in } A \\ - E \text{ quadrato} \end{array}\right\}$ æquari $D$ plano quæ æquatio est plani sub latere negati de quadrato. et argumentando ergo $\dfrac{E \text{ quadratum} + D \text{ plano}}{E}$ æquabitur $A$.

postremo per modos compositos, et excogitanda ab artifice et tentanda, quæ suo fini magis inservire conjicit, figmenta.

Quamcumq; autem speciem induit $A$, æquatio Juxta eam transformabitur et nova de $E$ ordinabitur. si quæ de $A$ enunciantur in æquatione proposita, eadem enuncientur de nova quam $A$ induit specie.

Alteretur enim $A$ per additionem, fingendo $A + B$ esse $E$ ut supra unde $E - B$ fit $A$. Cum igitur par potestas eliciatur ab $E - B$, quæ proponebatur ab $A$, et similes quoq; parodici gradus, in quos ducantur Invariandæ coefficientes, ad efficienda eadem adfectionum homogenea; omnino facta huiusmodi, æquabuntur proposito homogeneo. \uncertain{Expediet} igitur secundum artis præcepta, et ita demum æquatio de $E$ ordinetur. Jam transmutata erit in novam æquationem de $E$, id est ipsa $A$ producta \uncertain{incremento} $B$ enunciandam quod faciendum erat.

quod in reliquis quibuscumq; alterandi modis, quibus semper $A$ valor exprimitur locum habere conspicuum est.

\subsection*{De transmutatione Invariata radice}

quod ad secundam permutationem attinet, Invariata radice æquationem transmutare, est gradum æquationis deprimere vel attollere.
\note{La page finit ici, au trois quarts ; la suite est attendue au recto suivant.}

\page{20}
\note{Recto d'un nouveau feuillet, main du texte ; en haut à droite, à l'encre, le chiffre « 8 ». Le texte continue sans rupture la vue 19 (« … deprimere vel attollere. »). Au bord gauche de la vue paraît la marge du feuillet précédent, avec les fins de lignes de l'addition de la vue 18, et, dans la marge de gauche de ce feuillet-ci, trois courtes notes d'une petite main, reliées par des croix « + » à des croix posées dans le texte : ce sont les fragments aperçus en bord de la vue 19. À droite, au-delà du bord du feuillet, dépassent d'autres feuillets du cahier : l'un porte en haut un « 11 », un autre plus bas un « 9 », avec des fins de mots (« …p ad », « et », « …ctti », « …onis », « atq; »…) ; ils ne sont pas transcrits. Au pied, sous le bord du feuillet, paraissent comme aux vues 18 les dernières lignes d'un feuillet plus long : « …reductione viam, … altarum, et adfectarum. Æquationis tamen modo constitutarum spectata proprietas maxime Juvat. »}

Climacticus autem sive adscensus sive descensus regulariter fit vel irregulariter.

Climacticus regularis ascensus fit, cum utraq; propositæ æquationis pars, distributa ordinate, vel sine ea artificis industria ducitur quadratice, cubice et ulterius climactice.

descensus contra cum dividitur sub quadratice, subcubice et depressius subclimactice, adiectione, ne ad id supplemento dato vel inquirendo. omnia autem post ductionem, divisionemve congrue ordinantur.
\note{« ne ad id » est lu tel quel ; la phrase parallèle plus bas porte « adiecto nempe ad id supplemento ».}

Irregularis autem adscensus fit, cum omnia quibus proposita æquatio constat, singularia homogenea ducuntur in eundem parodicum gradum, sive purum\marginal{+ in exemplari \uncertain{excripto} / vel potius \uncertain{adscriptum}}, sive adfectum a datis coefficientibus magnitudinibus, vel etiam eas adficientium, et ducta congruam interpretationem arripiunt, et ordinationem.
\note{Deux petites croix « + » sont posées au-dessus de « purum » et de « eas » ; la note marginale correspondante, d'une petite main pâle, n'est lue qu'en partie.}

descensus contra cum omnia singularia illa homogenea adplicantur, sive ad radicem quæsititiam, sive radicis gradum potestate inferiorem, vel puro, vel cum adficientibus adfectione datis coefficientibus magnitudinibus, et divisa adiecto nempe ad id supplemento dato vel inquirendo congruam interpretationem arripiunt, et ordinationem.

demonstratio vero in quacumq; artificiosa transmutandi forma statim evidens fit, quoniam quæ æqualia sunt longitudine, æqualia quoq; sunt similis potestate, et contra. neq; communis divisor, vel multiplicator æqualitatem Immutat vel rationem.

Atq; hæc ut a nobis generaliter proposita sunt, ita specialibus inde regulæ præceptis, et exemplis concedendum est, quæ diffunduntur passim in commentationes locos, prout eadem artificij transmutatorij exigunt specimina.

\subsection*{Singularia de plasmate. Cap. VIII}

plasma Inesse æquationibus intelligitur cum deducuntur potestates adfectæ a puris, vel minus adfectis æqualitatis aut etiam altiores a depressioribus.

Itaq; plasma potest institui per omnes transmutationum modos, præter divisionem et climacticum descensum.

finis plastices ac præcipuus illius usus est, ut agnitæ æquationes plasmaticæ, ad simpliciores\marginal{+ vel in} a quibus deductæ sunt reducantur resolui, si quo id liceat compendio
\note{Une croix « + » au-dessus de « ad simpliciores » renvoie à la note marginale « + vel in ». « liceat » est écrit sur un premier mot. Le texte s'arrête ici, sans ponctuation, au bas de la partie écrite du feuillet ; la phrase est inachevée sur cette vue et continue vraisemblablement au verso (vue 21, hors de ce lot).}

\end{document}
